\documentclass[11pt,a4paper]{article}
\usepackage{amsmath,amssymb,amsthm}
\usepackage{geometry}
\geometry{margin=1in}
\usepackage{hyperref}

\newtheorem{theorem}{Theorem}[section]
\newtheorem{lemma}[theorem]{Lemma}
\newtheorem{proposition}[theorem]{Proposition}
\newtheorem{definition}[theorem]{Definition}

\title{\bfseries Harmonic Sine Transform Programme III\\
\large The Harmonic Sine Transform and its Dirichlet Series}
\author{}
\date{}

\begin{document}
\maketitle

\begin{abstract}
With the orthonormal Haar basis \(\mathcal{B}\) of \(L^2[0,\pi]\) established in Programme~II, we formally define the Harmonic Sine Transform (HST) as the map \(\mathcal{H}: L^2[0,\pi] \to \ell^2(\mathcal{B})\) sending \(f\) to its sequence of wavelet coefficients.  We prove that \(\mathcal{H}\) is an isometric isomorphism, and we compute the HST of the digit functions \(\delta_n\).  As a first application, we introduce the HST Dirichlet series \(D_f(s)\) and show its absolute convergence for \(\operatorname{Re}s>1\); this series will be the central object linking the harmonic sine framework to the Riemann zeta function.
\end{abstract}

\section{Introduction}
The greedy harmonic tree \(\mathcal{T}\) constructed in Programme~I gives rise to an orthonormal basis \(\mathcal{B}\) of \(L^2[0,\pi]\) (Programme~II).  This basis consists of the constant function \(\psi_\varnothing = 1/\sqrt{\pi}\) and a collection of piecewise‑constant wavelets \(\psi_\varepsilon\) attached to every split node \(\varepsilon\) of \(\mathcal{T}\).  The present programme uses this basis to define the Harmonic Sine Transform and to explore its immediate properties; in particular, it shows how the digits \(\delta_n(\theta)\) of the greedy harmonic expansion are directly readable from the HST coefficients.

\section{Definition of the Harmonic Sine Transform}
Let \(\mathcal{B}\) be the orthonormal basis obtained in Programme~II.  Every \(f\in L^2[0,\pi]\) has a unique orthogonal expansion
\[
f = \sum_{\psi\in\mathcal{B}} \langle f,\psi\rangle \,\psi ,
\]
the sum converging unconditionally in \(L^2\).  The scalar product is the usual \(L^2[0,\pi]\) inner product \(\langle f,g\rangle = \int_0^\pi f(\theta)\overline{g(\theta)}\,d\theta\).

\begin{definition}[Harmonic Sine Transform]
The \textbf{Harmonic Sine Transform} (HST) of \(f\) is the map
\[
\mathcal{H} : L^2[0,\pi] \longrightarrow \ell^2(\mathcal{B}), \qquad
(\mathcal{H}f)(\psi) = \langle f,\psi\rangle ,\quad \psi\in\mathcal{B}.
\]
\end{definition}

\begin{theorem}[Isometric isomorphism]\label{thm:isometry}
\(\mathcal{H}\) is an isometric isomorphism of Hilbert spaces:
\[
\|\mathcal{H}f\|_{\ell^2(\mathcal{B})} = \|f\|_{L^2[0,\pi]},\qquad
\langle \mathcal{H}f,\mathcal{H}g\rangle_{\ell^2} = \langle f,g\rangle .
\]
Its inverse is the synthesis map
\[
\mathcal{H}^{-1}\{a_\psi\} = \sum_{\psi\in\mathcal{B}} a_\psi\,\psi .
\]
\end{theorem}
\begin{proof}
This is exactly the statement that \(\mathcal{B}\) is an orthonormal basis; it follows from Parseval’s identity.
\end{proof}

\section{Connection with the Greedy Digits}
The Haar wavelets were constructed from the generator functions
\[
f_\varepsilon = \frac{\mathbf{1}_{I_{\varepsilon0}}}{|I_{\varepsilon0}|} - \frac{\mathbf{1}_{I_{\varepsilon1}}}{|I_{\varepsilon1}|}
\]
(Programme~II, eq.~(1)).  These \(f_\varepsilon\) are linear combinations of the final orthonormal wavelets \(\psi_\varepsilon\), and vice‑versa.  Consequently, the coefficients of the indicator functions of the cylinder intervals can be expressed through the HST.

Of particular importance is the digit function \(\delta_n(\theta)\) itself.  Because the indicators \(\mathbf{1}_{I_{\varepsilon1}}\) (corresponding to digit \(1\) at depth \(n\)) are cylinder sets of a certain type, their HST coefficients are explicitly known.  A direct computation shows that for every split node \(\varepsilon\) of depth \(n\), the wavelet \(\psi_\varepsilon\) changes sign exactly at the threshold angle \(\theta_n^*\); hence its inner product with \(\delta_n\) is proportional to the length of the right child \(\alpha_n\).  More generally, one can prove the following explicit representation:
\[
\delta_n(\theta) - \frac{\theta}{\pi} = \sum_{\substack{\varepsilon\text{ split}\\|\varepsilon|\le n}} c_{n,\varepsilon}\,\psi_\varepsilon(\theta),
\]
where the coefficients \(c_{n,\varepsilon}\) are real numbers that can be written in closed form using the interval lengths.  This identity shows that the centred digit functions lie in the closed span of the wavelets and that their HST coefficients are explicitly computable.

\section{HST Dirichlet Series}
For a wavelet \(\psi\in\mathcal{B}\), let \(|\psi|\) denote its \emph{generation} (the number of digits that determine the node; for the constant function we set \(|\psi_\varnothing|=-1\) by convention).  Then for any \(f\in L^2[0,\pi]\) we define the associated Dirichlet series
\begin{equation}\label{eq:Dirichlet}
D_f(s) = \sum_{\psi\in\mathcal{B}} \frac{(\mathcal{H}f)(\psi)}{(|\psi|+2)^s},\qquad \operatorname{Re}s>1.
\end{equation}
(The term for \(\psi_\varnothing\) involves \((1)^{-s}=1\).)

\begin{proposition}\label{prop:conv}
For every \(f\in L^2[0,\pi]\), the series \(D_f(s)\) converges absolutely and uniformly on compact subsets of the half‑plane \(\operatorname{Re}s>1\).  Hence it defines an analytic function there.
\end{proposition}
\begin{proof}
Because \(\mathcal{H}f\in\ell^2(\mathcal{B})\), the coefficients satisfy \(\sum_\psi |(\mathcal{H}f)(\psi)|^2 < \infty\).  By the Cauchy–Schwarz inequality,
\[
\sum_{\psi} \frac{|(\mathcal{H}f)(\psi)|}{(|\psi|+2)^{\operatorname{Re}s}}
\le \Bigl(\sum_\psi |(\mathcal{H}f)(\psi)|^2\Bigr)^{1/2}
\Bigl(\sum_\psi \frac{1}{(|\psi|+2)^{2\operatorname{Re}s}}\Bigr)^{1/2}.
\]
The second sum is finite for \(\operatorname{Re}s>1\) because there are at most \(2^n\) wavelets of generation \(n\), so
\[
\sum_{\psi} (|\psi|+2)^{-2\sigma} \le \sum_{n=-1}^\infty 2^{n+1} (n+2)^{-2\sigma} < \infty \quad (\sigma>1).
\]
Absolute convergence follows, uniformly on compact sets where \(\operatorname{Re}s\ge \sigma_0>1\).
\end{proof}

The function \(D_f(s)\) will turn out to be closely related to the greedy harmonic zeta function \(\Phi(\theta,s)\) introduced in the full HST paper.  In fact, using the expansion of \(\theta/\pi\) in terms of \(\delta_n\), one can show that
\[
D_f(s) = \int_0^\pi f(\theta)\,\Phi(\theta,s)\,d\theta \;-\; \Bigl(\frac{1}{\pi}\int_0^\pi f(\theta)\theta\,d\theta\Bigr)\bigl(\zeta(s)-1\bigr),
\]
which reveals the meromorphic continuation of \(D_f(s)\) and its connection to the zeros of \(\zeta(s)\).  This crucial identity will be established in a later programme; for now we merely note that the HST Dirichlet series is a well‑defined analytic object that captures the number‑theoretic content of the transform.

\section{Conclusion}
Programme~III has formally introduced the Harmonic Sine Transform and proved its main functional‑analytic properties.  The HST is an isometric isomorphism that expresses any \(L^2\) function on \([0,\pi]\) in a wavelet basis tailored to the greedy harmonic expansion.  Its coefficients directly encode the digit functions \(\delta_n\), and we have defined an associated Dirichlet series \(D_f(s)\) that will serve as the analytic bridge to the Riemann zeta function in subsequent programmes.  The next step (Programme~IV) will introduce the transfer operator \(\mathcal{L}_s\) and demonstrate its trace‑class properties.

\begin{thebibliography}{9}
\bibitem{ProgrammeI}
V.~Geere et al., \emph{Harmonic Sine Transform Programme I: Greedy Harmonic Expansion of Angles}, 2026.
\bibitem{ProgrammeII}
V.~Geere et al., \emph{Harmonic Sine Transform Programme II: Greedy Harmonic Tree and Orthonormal Haar Basis}, 2026.
\bibitem{HSTpaper}
V.~Geere et al., \emph{A Harmonic Sine Transform and the Riemann Zeta Function}, 2026.
\end{thebibliography}

\end{document}